\documentclass[12pt]{article}

\usepackage{sbc-template}
\usepackage[utf8]{inputenc}
\usepackage[brazil]{babel}
\usepackage{graphicx}
\usepackage[hyphens]{url}
\usepackage{float}
\usepackage{booktabs}
\usepackage{multirow}
\usepackage{tabularx}
\usepackage{array}

% colunas que quebram linha, alinhadas à esquerda (tabelas 2, 3 e 4)
\newcolumntype{L}{>{\raggedright\arraybackslash}X}
\newcommand{\cod}[1]{\texttt{#1}}

\sloppy

\title{DynamoDB vs SGBD Relacional: Um Estudo de Caso Comparativo em Big Data}

\author{Gabriel V. Cascaes\inst{1}, Letícia A. de O. Burlinski\inst{1},
  Vinícius E. D. da Silva\inst{1}, Uriel D. de Oliveira\inst{1}}

\address{Setor de Educação Profissional e Tecnológica\\
  Universidade Federal do Paraná (UFPR) -- Curitiba, PR -- Brasil
%  \email{PREENCHER@ufpr.br, PREENCHER@ufpr.br, PREENCHER@ufpr.br, PREENCHER@ufpr.br}
}

\begin{document}

\maketitle

\begin{abstract}
The growth in the volume, velocity and variety of data has challenged
traditional relational Database Management Systems (DBMSs) and driven the
adoption of NoSQL solutions. This paper presents a comparative case study
between Amazon DynamoDB, a managed NoSQL database based on the key-value and
document models, and MySQL, a relational DBMS, in storing and handling Big
Data. Using the same dataset and the same operations on both technologies, we
evaluate data modeling, performance, scalability and implementation
complexity. Preliminary results show that MySQL was faster in all evaluated
operations, by 10 to 65 times in key-based operations and by up to 500 times
in queries that require a full table scan in DynamoDB, highlighting
DynamoDB's dependence on previously modeled access patterns.
\end{abstract}

\begin{resumo}
O crescimento do volume, da velocidade e da variedade dos dados tem desafiado
os Sistemas Gerenciadores de Banco de Dados (SGBDs) relacionais tradicionais e
impulsionado a adoção de soluções NoSQL. Este artigo apresenta um estudo de
caso comparativo entre o Amazon DynamoDB, banco de dados NoSQL gerenciado
baseado nos modelos chave-valor e de documentos, e o MySQL, SGBD relacional,
no armazenamento e na manipulação de Big Data. A partir do mesmo conjunto de
dados e das mesmas operações em ambas as tecnologias, são avaliados aspectos
como modelagem de dados, desempenho, escalabilidade e complexidade de
implementação. Os resultados preliminares mostram que o MySQL foi mais rápido
em todas as operações avaliadas, com diferença de 10 a 65 vezes nas operações
por chave e de até 500 vezes nas consultas que, no DynamoDB, exigem varredura
completa da tabela, evidenciando a dependência do DynamoDB em padrões de
acesso previamente modelados.
\end{resumo}

\section{Introdução}

O crescimento acelerado do volume de dados gerados por aplicações web,
dispositivos móveis e plataformas de comércio eletrônico tem imposto
novos desafios ao armazenamento e ao processamento da informação. Esse
cenário, comumente chamado de \textit{Big Data}, é caracterizado não apenas
pelo grande volume de dados, mas também pela alta velocidade com que são
produzidos e consumidos e pela variedade de formatos em que se apresentam.
Nesse contexto, as organizações passaram a demandar soluções capazes de
escalar horizontalmente, manter baixa latência sob cargas elevadas e oferecer
alta disponibilidade a todo momento.

Por décadas, os Sistemas Gerenciadores de Banco de Dados (SGBDs) relacionais
foram a principal escolha para a persistência de dados, graças ao modelo
relacional, à linguagem SQL e às garantias transacionais
ACID. Contudo, à medida que as aplicações passaram a operar em escala global,
algumas características desse modelo, como o esquema rígido, a dependência de
junções (JOINs) e a dificuldade de distribuir dados entre múltiplos servidores,
se tornaram limitantes em determinados cenários. Esse
contexto impulsionou o surgimento dos bancos de dados NoSQL, que relaxam
parte dessas garantias em favor de escalabilidade, flexibilidade e desempenho.

Diante disso, este trabalho tem como objetivo avaliar comparativamente o uso
do Amazon DynamoDB e do MySQL no armazenamento e na manipulação de dados em
um cenário de \textit{Big Data}. Para isso, foi desenvolvido um estudo de
caso, no qual o
mesmo conjunto de dados e as mesmas operações foram executados em ambas as
tecnologias. A análise busca identificar as diferenças, as vantagens e as
limitações de cada abordagem, considerando aspectos como modelagem de dados,
desempenho, escalabilidade e complexidade de implementação.


\section{Tecnologia Utilizada}
\subsection{Amazon DynamoDB}

O Amazon DynamoDB é um serviço de banco de dados NoSQL totalmente gerenciado,
oferecido pela Amazon Web Services (AWS) desde 2012, que suporta os modelos de
dados chave-valor e de documentos. Por ser um serviço \textit{serverless}, o
usuário não administra servidores, sistema operacional, replicação ou
\textit{backups}: a AWS se encarrega do particionamento automático dos dados,
da replicação entre múltiplas zonas de disponibilidade e da manutenção da
infraestrutura. Essa tecnologia tem origem no Dynamo, sistema de armazenamento chave-valor
desenvolvido internamente pela Amazon \cite{decandia2007dynamo}.

\subsection{Modelagem de Dados}

No DynamoDB, os dados são organizados em \textbf{tabelas}, compostas por
\textbf{itens}, que por sua vez são compostos por \textbf{atributos}. Um item
é análogo a uma linha de uma tabela relacional, porém sem esquema fixo: além
da chave primária, cada item pode possuir um conjunto diferente de atributos.
Os atributos podem ser de tipos escalares (\textit{string}, número, binário,
booleano e nulo), de conjunto ou de documento (listas e mapas, que permitem
estruturas aninhadas semelhantes a JSON). Cada item está limitado a 400~KB
\cite{awsdynamodbdocs}.

A chave primária é o elemento central da modelagem e pode ser de dois tipos:

\begin{itemize}
  \setlength{\itemsep}{0pt}
  \item \textbf{Simples}: composta apenas pela chave de partição
  (\textit{partition key}), cujo valor passa por uma função \textit{hash}
  que determina em qual partição física o item será armazenado;
  \item \textbf{Composta}: formada pela chave de partição e por uma chave de
  ordenação (\textit{sort key}). Itens com a mesma chave de partição são
  armazenados juntos e ordenados pela chave de ordenação, o que permite
  consultas por intervalo (por exemplo, todos os pedidos de um cliente em
  determinado período).
\end{itemize}

Para consultar os dados por atributos que não fazem parte da chave primária,
o DynamoDB oferece índices secundários: o \textbf{índice secundário local}
(LSI), que mantém a mesma chave de partição e utiliza uma chave de ordenação
alternativa, e o \textbf{índice secundário global} (GSI), que permite
definir chaves de partição e de ordenação completamente diferentes das da
tabela.

A principal diferença em relação ao modelo relacional está na abordagem de
modelagem. Enquanto no modelo relacional os dados são normalizados com base
nas entidades e em seus relacionamentos, e as consultas são definidas
posteriormente com o uso de junções, no DynamoDB a modelagem parte dos
\textbf{padrões de acesso} da aplicação: primeiro definem-se quais consultas
serão realizadas e, a partir delas, as chaves e os índices são criados. Como o DynamoDB não suporta junções, é comum a
desnormalização dos dados e a adoção do \textit{single-table design}, em que
entidades diferentes e relacionadas são armazenadas em uma mesma tabela,
distinguidas por prefixos nas chaves, de modo que uma única operação de
consulta recupere todos os dados necessários.

Essa abordagem garante eficiência nas consultas previstas, mas reduz a
flexibilidade: consultas não planejadas exigem a criação de novos índices ou
a varredura completa da tabela (operação \textit{Scan}), que é custosa e
lenta em grandes volumes de dados.

\subsection{Uso no Mercado e Comparação com o Modelo Relacional}

O DynamoDB é amplamente utilizado em aplicações que exigem alta
escalabilidade e baixa latência. A própria Amazon o utiliza em sistemas
críticos de suas operações: durante o Prime Day de 2026, o serviço processou
mais de 59 trilhões de requisições, com pico de 192 milhões de requisições
por segundo, mantendo tempos de resposta de poucos milissegundos
\cite{aws2026primeday}. Fora da Amazon, o serviço é adotado por empresas
como Disney+, Nike, Experian e GE Aerospace, em setores como varejo,
\textit{streaming} de mídia, jogos, publicidade digital e serviços
financeiros.

De forma geral, o DynamoDB tende a ser mais vantajoso que um SGBD relacional
quando:

\begin{itemize}
  \setlength{\itemsep}{0pt}
  \item a aplicação precisa escalar horizontalmente para volumes muito altos
  de leitura e escrita, com picos imprevisíveis de tráfego;
  \item os padrões de acesso são conhecidos e estáveis, baseados
  principalmente em buscas por chave;
  \item a latência baixa e previsível é um requisito crítico;
  \item os dados têm estrutura flexível ou variável entre registros;
  \item a equipe deseja reduzir o esforço de administração de
  infraestrutura de banco de dados.
\end{itemize}

\section{Estudo de Caso}

O estudo de caso utiliza o \textit{Brazilian E-Commerce Public Dataset by
Olist} \cite{olist2018}, conjunto público e anonimizado de cerca de 100 mil
pedidos realizados entre 2016 e 2018 em \textit{marketplaces} brasileiros. Os
dados estão distribuídos em nove arquivos CSV, dos quais oito foram
utilizados: clientes (99.441 registros), produtos (32.951), vendedores
(3.095), pedidos (99.441), itens dos pedidos (112.650), pagamentos (103.886),
avaliações (99.224) e geolocalização (cerca de 1 milhão de registros). A
Figura~\ref{fig:esquema} apresenta os arquivos e seus relacionamentos, e o Anexo~\ref{anexo:modelagem}
detalha os principais atributos.


\begin{figure}[htbp]
    \centering
    \includegraphics[width=0.5\linewidth]{image.png}
    \caption{Esquema de dados do conjunto Olist. Fonte: \cite{olist2018}}
    \label{fig:esquema}
\end{figure}

\subsection{Questões Comparativas}

As mesmas questões foram executadas nos dois bancos. Elas cobrem operações em
que o DynamoDB tende a ser vantajoso (acessos por chave) e operações em que
tende a ser limitado (buscas textuais, junções e agregações). Entre
parênteses, o resultado esperado:

\begin{enumerate}
  \setlength{\itemsep}{0pt}
  \item \textbf{Acesso por chave:} dados de um pedido, dado o seu ID (o
  pedido, seus itens, pagamentos e avaliação).
  \item \textbf{Inserção:} criação de um pedido com itens e pagamento
  (registros gravados de forma atômica).
  \item \textbf{Atualização:} alteração do status de um pedido (um registro
  alterado).
  \item \textbf{Carga:} tempo para carregar o conjunto completo (o mesmo
  número de registros nos dois bancos).
  \item \textbf{Filtro por intervalo:} pedidos entregues em janeiro de 2018
  (um registro por pedido).
  \item \textbf{Agregação em uma tabela:} valor médio, mínimo e máximo pago em
  cada forma de pagamento, da maior média para a menor (uma linha por forma).
  \item \textbf{Busca de texto em uma tabela:} avaliações que mencionam
  ``atraso'' no comentário (as avaliações com o termo).
  \item \textbf{Busca de texto em duas tabelas:} pedidos de clientes de
  cidades cujo nome contém ``sao'' (um registro por pedido).
  \item \textbf{Agregação em três tabelas:} preço médio, mínimo e máximo dos
  itens por UF do cliente, da maior média para a menor (uma linha por UF).
\end{enumerate}

\section{Resultados e Avaliações}
\subsection{Configuração do Ambiente}
Os experimentos foram executados em uma máquina com processador Intel Core i5-8250U (até 3,40~GHz), 16~GB de memória RAM
e sistema operacional Ubuntu 22.04. O SGBD relacional
utilizado foi o MySQL, executado localmente. O Amazon DynamoDB foi utilizado na região us-east-2.


\subsection{Modelagem dos Dados}

Conforme a Seção~2.2, no MySQL o esquema foi derivado das entidades e, no
DynamoDB, dos padrões de acesso da Seção~3.1. Chaves, índices e operações de
cada questão estão detalhados no Anexo~\ref{anexo:modelagem}.

\textbf{MySQL.} Esquema normalizado com uma tabela por arquivo do Olist, chaves
estrangeiras entre pedidos, itens, pagamentos, avaliações e clientes,
identificadores em \texttt{CHAR(32)}, datas em \texttt{DATETIME} e valores em
\texttt{DECIMAL(10,2)}. Além dos índices das chaves, foram criados um índice em
\texttt{order\_delivered\_customer\_date}, para o filtro por data de entrega, e
um índice composto (\texttt{payment\_type}, \texttt{payment\_value}), para a
agregação por forma de pagamento.

\textbf{DynamoDB.} Adotou-se o \textit{single-table design}: uma única tabela,
\texttt{Olist}, no modo sob demanda, com chave composta \texttt{PK}/\texttt{SK}
e prefixos que identificam cada entidade. O pedido é o agregado central: o
item de metadados (\texttt{SK = METADATA}), os itens (\texttt{ITEM\#001},
\ldots), os pagamentos e a avaliação compartilham a chave
\texttt{PK = ORDER\#<order\_id>}, de modo que uma única \texttt{Query}
recupera o pedido completo. A cidade e a UF do cliente foram desnormalizadas no
item do pedido, e a categoria do produto, em cada item. Dois índices
secundários globais, com projeção \texttt{KEYS\_ONLY}, atendem às consultas que
não partem do pedido: o GSI1, esparso, agrupa os pedidos entregues por mês
(\texttt{DELIVERED\#<AAAA-MM>}), ordenados pela data de entrega; o GSI2 agrupa
os pagamentos por forma de pagamento (\texttt{PAYTYPE\#<tipo>}), ordenados pelo
valor pago. As buscas de texto e a agregação por UF não têm índice
correspondente e exigem \textit{Scan}.

\subsection{Resultados Preliminares}

Os resultados a seguir são preliminares, obtidos em uma máquina com
Windows~11, com o MySQL local em contêiner Docker e o DynamoDB na região
\texttt{us-east-2}, sem a geolocalização; os testes definitivos serão
repetidos na máquina da Seção~4.1, com o conjunto completo. Em todas as
questões, os dois bancos retornaram o mesmo número de linhas, o que confirma a
equivalência das implementações.

\textbf{Carga.} O MySQL carregou cerca de 550 mil registros em 10,3~s com
\texttt{LOAD DATA LOCAL INFILE}, enquanto o DynamoDB levou 540,2~s, cerca de
52 vezes mais, com uma taxa aproximada de 1.000 itens por segundo
(Tabela~\ref{tab:carga}, no Anexo~\ref{anexo:resultados}). A diferença decorre
do limite de 25 itens por requisição \texttt{BatchWriteItem} e da latência de
rede entre o cliente e a região da AWS, ausente no MySQL local.

\textbf{Consultas.} A Tabela~\ref{tab:consultas} apresenta a média de três
execuções de cada questão e a operação usada no DynamoDB; o número de linhas
retornadas está no Anexo~\ref{anexo:resultados}.

\begin{table}[htbp]
\centering
\caption{Tempo médio das consultas (ms, média de 3 execuções)}
\label{tab:consultas}
\footnotesize
\begin{tabular}{|l|l|r|r|r|}
\hline
\textbf{Questão} & \textbf{Operação no DynamoDB} & \textbf{MySQL} & \textbf{DynamoDB} & \textbf{Razão} \\ \hline
Acesso por chave             & \texttt{Query}             & 9,25   & 599,70    & 65$\times$  \\ \hline
Inserção de pedido           & \texttt{TransactWriteItems}& 17,52  & 174,83    & 10$\times$  \\ \hline
Atualização de status        & \texttt{UpdateItem}        & 11,26  & 161,45    & 14$\times$  \\ \hline
Entregues em jan/2018        & \texttt{Query} (GSI1)      & 43,71  & 1.541,95  & 35$\times$  \\ \hline
Estatísticas por pagamento   & \texttt{Query} (GSI2)      & 45,19  & 22.441,57 & 497$\times$ \\ \hline
Busca de texto (1 tabela)    & \texttt{Scan}              & 69,08  & 23.044,77 & 334$\times$ \\ \hline
Busca de texto (2 tabelas)   & \texttt{Scan}              & 295,94 & 23.674,01 & 80$\times$  \\ \hline
Agregação por UF (3 tabelas) & \texttt{Scan}              & 775,44 & 40.423,14 & 52$\times$  \\ \hline
\end{tabular}
\end{table}


As operações por chave foram resolvidas pelo DynamoDB com uma única
requisição cada, em 160 a 600~ms, tempo dominado pela ida e volta até a região
da AWS, e não pelo processamento no banco; o valor maior do acesso por chave,
primeira operação da bateria, provavelmente inclui a abertura da conexão. A
consulta de entregas de janeiro de 2018 mostra o efeito do GSI1: o DynamoDB
leu apenas as 6.597 chaves do mês. Já as estatísticas por forma de pagamento,
mesmo com o GSI2, exigiram ler as cerca de 104 mil entradas do índice e
calcular média, mínimo e máximo na aplicação, pois o DynamoDB não oferece
funções de agregação.

As buscas de texto e a agregação por UF dependem do \textit{Scan}, que lê a
tabela inteira (cerca de 550 mil itens) antes de aplicar o filtro: seu tempo,
entre 23 e 40~s, é proporcional ao tamanho da tabela, e não ao resultado
(305 e 21.092 linhas levaram o mesmo tempo nas duas buscas). A
desnormalização eliminou as junções, mas não a varredura. No MySQL, as mesmas
consultas levaram menos de 1~s com \texttt{LIKE}, \texttt{JOIN} e
\texttt{GROUP BY}. Ressalta-se que o experimento compara um banco local com um
serviço remoto, com um único cliente sequencial, cenário que não explora a
escalabilidade horizontal e a latência estável sob alta concorrência do
DynamoDB e penaliza cada requisição com a latência de rede.

\section{Conclusão}

Os resultados preliminares confirmam a hipótese do estudo: no DynamoDB, o
desempenho depende de a consulta ter sido prevista na modelagem. Operações por
chave e consultas atendidas por índices secundários foram resolvidas com poucas
requisições, enquanto buscas textuais, junções e agregações exigiram a leitura
completa da tabela e processamento na aplicação, ficando de 50 a 500 vezes mais
lentas que no MySQL. Para cargas analíticas e consultas não planejadas, o SGBD
relacional mostrou-se mais simples e eficiente; o DynamoDB é recomendado quando
os padrões de acesso são conhecidos, baseados em chave, e a aplicação exige
escalabilidade e disponibilidade sem administração de infraestrutura.

\bibliographystyle{sbc}
\bibliography{sbc-template}

\newpage
\appendix
\section{Guia de Uso do Amazon DynamoDB na AWS}
\label{anexo:guia}

Este anexo é um guia prático para criar, carregar, consultar, monitorar e
remover tabelas do Amazon DynamoDB em uma conta da AWS, pelo console web, pela
linha de comando (AWS CLI) e por programas em Python (biblioteca
\texttt{boto3}). Os exemplos usam a tabela \texttt{Olist} deste trabalho.

\subsection{Acesso à Conta e Credenciais}

\textbf{Console e região.} O DynamoDB é acessado pelo console da AWS
(\texttt{console.aws.amazon.com}), pesquisando ``DynamoDB'' na barra de busca.
As tabelas pertencem a uma região: uma tabela criada em
\texttt{us-east-2} (Ohio) não aparece quando o console está em outra região.
A região atual é exibida no canto superior direito do console e no endereço
da página (\texttt{region=us-east-2}).

\textbf{Permissões.} O acesso é controlado pelo IAM (\textit{Identity and
Access Management}). Recomenda-se não usar o usuário raiz da conta, e sim um
usuário ou papel com as permissões necessárias, por exemplo a política
gerenciada \texttt{AmazonDynamoDBFullAccess} ou uma política restrita às
ações usadas: \texttt{CreateTable}, \texttt{DescribeTable},
\texttt{ListTables}, \texttt{DeleteTable}, \texttt{PutItem},
\texttt{BatchWriteItem}, \texttt{UpdateItem}, \texttt{DeleteItem},
\texttt{GetItem}, \texttt{Query} e \texttt{Scan}.

\textbf{AWS CLI.} Após instalar a AWS CLI versão 2, há duas formas de
configurar as credenciais:

\begin{itemize}
  \item \texttt{aws login}: abre o navegador para autenticação com a conta da
  AWS e gera credenciais temporárias, renovadas automaticamente. Para que
  programas em Python usem esse login, é necessário o pacote
  \texttt{botocore[crt]};
  \item \texttt{aws configure}: grava a chave de acesso (\textit{Access Key ID}
  e \textit{Secret Access Key}) de um usuário IAM, gerada no console em IAM
  $\rightarrow$ Usuários $\rightarrow$ Credenciais de segurança.
\end{itemize}

As credenciais nunca devem ser escritas no código, em capturas de tela ou em
repositórios. Para testar a configuração:

{\footnotesize
\begin{verbatim}
aws dynamodb list-tables --region us-east-2
\end{verbatim}
}

\subsection{Criação de uma Tabela}

\textbf{Pelo console.} Em DynamoDB $\rightarrow$ Tabelas $\rightarrow$
\textit{Criar tabela}:

\begin{enumerate}
  \item informar o nome da tabela e a chave de partição (por exemplo,
  \texttt{PK}, do tipo \textit{String});
  \item opcionalmente, informar a chave de ordenação (\texttt{SK});
  \item em \textit{Configurações da tabela}, escolher \textit{Personalizar}
  para definir o modo de capacidade (sob demanda ou provisionado) e criar
  índices secundários globais;
  \item confirmar em \textit{Criar tabela} e aguardar o status
  \textit{Ativo}.
\end{enumerate}

\textbf{Pela AWS CLI.} O comando abaixo cria a tabela \texttt{Olist} com
chave composta, capacidade sob demanda e o índice GSI1:

{\footnotesize
\begin{verbatim}
aws dynamodb create-table --table-name Olist \
  --region us-east-2 --billing-mode PAY_PER_REQUEST \
  --attribute-definitions \
      AttributeName=PK,AttributeType=S \
      AttributeName=SK,AttributeType=S \
      AttributeName=GSI1PK,AttributeType=S \
      AttributeName=GSI1SK,AttributeType=S \
  --key-schema AttributeName=PK,KeyType=HASH \
               AttributeName=SK,KeyType=RANGE \
  --global-secondary-indexes '[{"IndexName":"GSI1",
      "KeySchema":[{"AttributeName":"GSI1PK","KeyType":"HASH"},
                   {"AttributeName":"GSI1SK","KeyType":"RANGE"}],
      "Projection":{"ProjectionType":"KEYS_ONLY"}}]'

aws dynamodb wait table-exists --table-name Olist \
  --region us-east-2
\end{verbatim}
}

Apenas os atributos usados em chaves e índices são declarados na criação; os
demais são definidos livremente em cada item.

\subsection{Boas Práticas de Modelagem das Chaves}

\begin{itemize}
  \item Partir das consultas da aplicação: cada padrão de acesso deve ser
  atendido por uma \texttt{Query} na tabela ou em um índice.
  \item Escolher chaves de partição com muitos valores distintos e acesso bem
  distribuído, evitando concentrar o tráfego em poucas chaves
  (``partições quentes'').
  \item Usar prefixos nos valores das chaves (\texttt{ORDER\#},
  \texttt{ITEM\#}) para guardar entidades diferentes na mesma tabela e
  preencher números com zeros à esquerda (\texttt{ITEM\#001}), pois a chave
  de ordenação do tipo texto é comparada caractere a caractere.
  \item Guardar datas no formato ISO~8601 (\texttt{AAAA-MM-DD hh:mm:ss}), cuja
  ordem alfabética coincide com a cronológica.
  \item Criar índices secundários globais para os acessos que não partem da
  chave principal; um índice esparso, preenchido só nos itens relevantes,
  ocupa menos espaço. A projeção \texttt{KEYS\_ONLY} reduz o tamanho do
  índice quando bastam as chaves.
  \item Respeitar os limites: 400~KB por item e, por padrão, até 20 índices
  secundários globais por tabela.
\end{itemize}

\subsection{Gravação de Dados}

\textbf{Pelo console.} Em \textit{Explorar itens} $\rightarrow$
\textit{Criar item}, os atributos podem ser preenchidos em formulário ou em
JSON. É útil para testes, não para cargas.

\textbf{Item a item (\texttt{PutItem}).} Grava ou substitui um item inteiro:

{\footnotesize
\begin{verbatim}
aws dynamodb put-item --table-name Olist --region us-east-2 \
  --item '{"PK":{"S":"ORDER#abc"},"SK":{"S":"METADATA"},
           "order_status":{"S":"delivered"}}'
\end{verbatim}
}

\textbf{Em massa (\texttt{BatchWriteItem}).} Grava até 25 itens por
requisição. Em Python, o \texttt{batch\_writer} agrupa os itens e reenvia
automaticamente os que a AWS não processar:

{\footnotesize
\begin{verbatim}
import boto3
dynamodb = boto3.resource("dynamodb", region_name="us-east-2")
tabela = dynamodb.Table("Olist")
with tabela.batch_writer() as lote:
    for item in itens:          # cada item é um dicionário Python
        lote.put_item(Item=item)
\end{verbatim}
}

Como cada requisição atravessa a rede, a carga é acelerada com vários envios
em paralelo. Valores numéricos devem ser do tipo \texttt{Decimal}, e não
\texttt{float}. Para volumes muito grandes, a AWS também permite criar uma
tabela a partir de arquivos armazenados no Amazon S3 (\textit{Importar do
S3}), sem consumir capacidade de escrita.

\subsection{Leitura de Dados}

\begin{itemize}
  \item \textbf{GetItem}: lê um item pela chave primária completa.
  \item \textbf{Query}: lê os itens de uma mesma chave de partição, na tabela
  ou em um índice, podendo filtrar a chave de ordenação por intervalo
  (\texttt{BETWEEN}) ou prefixo (\texttt{begins\_with}). É a forma eficiente
  de leitura.
  \item \textbf{Scan}: lê a tabela inteira. O filtro
  (\texttt{FilterExpression}) é aplicado depois da leitura, então o tempo e o
  custo dependem do tamanho da tabela, e não do resultado.
\end{itemize}

{\footnotesize
\begin{verbatim}
aws dynamodb query --table-name Olist --region us-east-2 \
  --key-condition-expression "PK = :pk" \
  --expression-attribute-values '{":pk":{"S":"ORDER#abc"}}'

aws dynamodb query --table-name Olist --region us-east-2 \
  --index-name GSI1 \
  --key-condition-expression "GSI1PK = :m" \
  --expression-attribute-values '{":m":{"S":"DELIVERED#2018-01"}}'
\end{verbatim}
}

Em Python, com paginação:

{\footnotesize
\begin{verbatim}
from boto3.dynamodb.conditions import Key
kw = dict(KeyConditionExpression=Key("PK").eq("ORDER#abc"))
resp = tabela.query(**kw)
itens = resp["Items"]
while "LastEvaluatedKey" in resp:     # resposta limitada a 1 MB
    resp = tabela.query(**kw,
        ExclusiveStartKey=resp["LastEvaluatedKey"])
    itens += resp["Items"]
\end{verbatim}
}

Cuidados na leitura:

\begin{itemize}
  \item cada resposta traz no máximo 1~MB; o restante deve ser buscado com
  \texttt{LastEvaluatedKey}, como acima;
  \item \texttt{ProjectionExpression} limita os atributos retornados e reduz
  o tráfego;
  \item atributos com nomes reservados (como \texttt{status} ou
  \texttt{name}) devem ser referenciados por nomes substitutos, em
  \texttt{ExpressionAttributeNames};
  \item o console oferece também o \textit{Editor PartiQL}, que aceita
  comandos semelhantes ao SQL (\texttt{SELECT}, \texttt{INSERT},
  \texttt{UPDATE}); um \texttt{SELECT} sem condição na chave é executado
  como \textit{Scan}.
\end{itemize}

\subsection{Atualizações e Transações}

\texttt{UpdateItem} altera atributos de um item sem regravá-lo inteiro:

{\footnotesize
\begin{verbatim}
tabela.update_item(
    Key={"PK": "ORDER#abc", "SK": "METADATA"},
    UpdateExpression="SET #s = :s",
    ExpressionAttributeNames={"#s": "order_status"},
    ExpressionAttributeValues={":s": "shipped"})
\end{verbatim}
}

\texttt{TransactWriteItems} grava, atualiza ou remove até 100 itens de forma
atômica: ou todas as operações são aplicadas, ou nenhuma. Neste trabalho,
foi usado para inserir um pedido junto com seus itens e pagamento.

\subsection{Monitoramento}

\begin{itemize}
  \item \textbf{Explorar itens}: mostra os itens da tabela e permite filtrar
  e consultar pelo console.
  \item \textbf{Visão geral}: exibe chaves, índices, modo de capacidade,
  número de itens e tamanho. Esses dois últimos valores são atualizados apenas
  a cada seis horas, aproximadamente; logo após uma carga, podem aparecer
  zerados. A opção \textit{Obter contagem de itens em tempo real} faz a
  contagem na hora, mas lê a tabela inteira.
  \item \textbf{Monitorar}: gráficos do Amazon CloudWatch com a capacidade de
  leitura e escrita consumida, a latência e as requisições limitadas
  (\textit{throttled}).
\end{itemize}

\subsection{Custos}

\begin{itemize}
  \item \textbf{Sob demanda} (\textit{on-demand}): cobra-se por requisição
  de leitura e escrita, sem planejamento de capacidade. Indicado para cargas
  variáveis e para experimentos como este.
  \item \textbf{Provisionado}: define-se a capacidade de leitura (RCU) e de
  escrita (WCU) por segundo, com escalonamento automático opcional. Uma RCU
  corresponde a uma leitura fortemente consistente de até 4~KB por segundo
  (ou duas eventualmente consistentes); uma WCU, a uma escrita de até 1~KB.
  \item Cobra-se também o armazenamento da tabela e dos índices.
  \item O \textit{Scan} é a operação mais cara, pois cada execução consome
  leitura proporcional ao tamanho da tabela.
  \item Recomenda-se configurar um alerta de orçamento no AWS Budgets e
  conferir as condições vigentes de créditos e nível gratuito da conta.
\end{itemize}

\subsection{Remoção dos Recursos}

Ao final dos testes, a tabela deve ser excluída, pelo console (selecionar a
tabela $\rightarrow$ \textit{Excluir}) ou pela CLI. A exclusão remove também
os índices e não pode ser desfeita:

{\footnotesize
\begin{verbatim}
aws dynamodb delete-table --table-name Olist --region us-east-2
\end{verbatim}
}

\subsection{Desenvolvimento Local}

Para desenvolver e testar sem custo, a AWS distribui o DynamoDB Local, que
pode ser executado em contêiner Docker e aceita as mesmas chamadas da API:

{\footnotesize
\begin{verbatim}
docker run -d -p 8000:8000 amazon/dynamodb-local
aws dynamodb list-tables --endpoint-url http://localhost:8000
\end{verbatim}
}

Em Python, basta informar \texttt{endpoint\_url="http://localhost:8000"} ao
criar o recurso. O DynamoDB Local não reproduz o desempenho nem a
escalabilidade do serviço; por isso, os resultados deste trabalho foram
medidos na AWS.

\subsection{Configuração Utilizada neste Trabalho}

Neste trabalho, foi criada uma única tabela, \texttt{Olist}, na região
\texttt{us-east-2}, no modo sob demanda, com chave composta
(\texttt{PK}, \texttt{SK}) e dois índices secundários globais com projeção
\texttt{KEYS\_ONLY}: GSI1 (pedidos entregues por mês) e GSI2 (pagamentos por
tipo). A carga e as consultas foram feitas pelo programa
\texttt{olist\_bench\_v2.py}, disponível em [LINK DO REPOSITÓRIO], usando as
operações descritas neste anexo.


\clearpage
\section{Detalhamento do Conjunto de Dados e da Modelagem}
\label{anexo:modelagem}

\subsection{Principais Atributos}

\begin{table}[htbp]
  \centering
  \caption{Principais atributos do conjunto de dados}
  \label{tab:atributos}
  \begin{tabularx}{\textwidth}{l l X}
    \toprule
    \textbf{Categoria do Atributo} & \textbf{Atributo Principal} & \textbf{Descrição / Papel no Negócio} \\
    \midrule
    \multirow{3}{*}{Identificadores}
      & \texttt{order\_id}    & Chave única que identifica cada pedido realizado. \\
      & \texttt{customer\_id} & Identificador da relação do cliente com a compra específica. \\
      & \texttt{product\_id}  & Identificador único do produto comercializado. \\
    \midrule
    \multirow{2}{*}{Dados Temporais}
      & \texttt{purchase\_timestamp}      & Data e hora exata em que o pedido foi efetuado pelo cliente. \\
      & \texttt{delivered\_customer\_date} & Data e hora real de entrega do pedido ao cliente final. \\
    \midrule
    \multirow{3}{*}{Valores Financeiros}
      & \texttt{price}         & Preço unitário do produto comercializado no item. \\
      & \texttt{freight\_value} & Valor do frete cobrado pelo transporte do item. \\
      & \texttt{payment\_value} & Valor monetário total pago através de uma determinada forma de pagamento. \\
    \bottomrule
  \end{tabularx}
\end{table}

\subsection{Modelo Relacional (MySQL)}

No MySQL, adotou-se um esquema normalizado que reproduz a estrutura original
do dataset (Figura~\ref{fig:esquema}), com uma tabela para cada arquivo:
\texttt{customers}, \texttt{orders}, \texttt{order\_items},
\texttt{order\_payments}, \texttt{order\_reviews}, \texttt{products},
\texttt{sellers} e \texttt{geolocation}. Os identificadores do Olist são
cadeias hexadecimais de 32 caracteres e foram armazenados como
\texttt{CHAR(32)}; os campos temporais, como \texttt{DATETIME}; e os
valores financeiros, como \texttt{DECIMAL(10,2)}, evitando erros de
arredondamento de tipos de ponto flutuante. As chaves e os índices
definidos são apresentados na Tabela~\ref{tab:mysql}.



\begin{table}[htbp]
\centering
\caption{Chaves e índices do modelo relacional}
\label{tab:mysql}
\footnotesize
\begin{tabularx}{\textwidth}{|>{\raggedright\arraybackslash}p{2.4cm}|L|L|L|}
\hline
\textbf{Tabela} & \textbf{Chave primária} & \textbf{Chaves estrangeiras} & \textbf{Índices adicionais} \\ \hline
\cod{orders} & \cod{order\_id} & \cod{customer\_id} & \cod{order\_delivered\_\allowbreak customer\_date} \\ \hline
\cod{order\_items} & (\cod{order\_id}, \cod{order\_item\_id}) & \cod{order\_id}, \cod{product\_id}, \cod{seller\_id} & -- \\ \hline
\cod{order\_\allowbreak payments} & (\cod{order\_id}, \cod{payment\_\allowbreak sequential}) & \cod{order\_id} & (\cod{payment\_type}, \cod{payment\_value}) \\ \hline
\cod{order\_\allowbreak reviews} & (\cod{review\_id}, \cod{order\_id}) & \cod{order\_id} & -- \\ \hline
\cod{customers} & \cod{customer\_id} & -- & \cod{customer\_zip\_\allowbreak code\_prefix} \\ \hline
\cod{products} & \cod{product\_id} & -- & -- \\ \hline
\cod{sellers} & \cod{seller\_id} & -- & \cod{seller\_zip\_\allowbreak code\_prefix} \\ \hline
\cod{geolocation} & \cod{geo\_id} (substituta, \cod{AUTO\_\allowbreak INCREMENT}) & -- & \cod{geolocation\_\allowbreak zip\_code\_prefix} \\ \hline
\end{tabularx}
\end{table}
 
Duas particularidades do dataset influenciaram o esquema. Primeiro, o
atributo \texttt{review\_id} não é único no arquivo de avaliações, o que
exigiu uma chave primária composta com \texttt{order\_id}. Segundo, a tabela
\texttt{geolocation} não possui chave natural, pois um mesmo prefixo de CEP
aparece em várias linhas com coordenadas diferentes; por isso, foi criada
uma chave substituta. Pelo mesmo motivo, e porque alguns prefixos de CEP de
clientes e vendedores não constam nessa tabela, a ligação por
\texttt{zip\_code\_prefix} foi mantida apenas por índices, sem restrição de
chave estrangeira.
 
Além dos índices criados automaticamente para chaves primárias e
estrangeiras, foram definidos dois índices voltados às consultas da
Seção~3.1: um sobre \texttt{order\_delivered\_customer\_date}, que permite
filtrar pedidos por intervalo de datas sem varrer a tabela, e um índice
composto (\texttt{payment\_type}, \texttt{payment\_value}) em
\texttt{order\_payments}, que cobre integralmente a consulta de agregação por
forma de pagamento.
 
\subsection{Modelo NoSQL (DynamoDB)}

 
No DynamoDB, adotou-se o \textit{single-table design}: todas as entidades
foram armazenadas em uma única tabela, \texttt{Olist}, com chave primária
composta por atributos genéricos \texttt{PK} (chave de partição) e
\texttt{SK} (chave de ordenação). O tipo de cada item é identificado por
prefixos nos valores das chaves, conforme a Tabela~\ref{tab:dynamo}.
 
\begin{table}[htbp]
\centering
\caption{Tipos de item na tabela única do DynamoDB}
\label{tab:dynamo}
\footnotesize
\begin{tabularx}{\textwidth}{|>{\raggedright\arraybackslash}p{2.7cm}|>{\raggedright\arraybackslash}p{2.9cm}|>{\raggedright\arraybackslash}p{2.6cm}|L|}
\hline
\textbf{Entidade} & \textbf{PK} & \textbf{SK} & \textbf{Principais atributos} \\ \hline
Pedido & \cod{ORDER\#\allowbreak<order\_id>} & \cod{METADATA} & status, datas, \cod{customer\_id}, cidade e UF do cliente, \cod{GSI1PK}, \cod{GSI1SK} \\ \hline
Item do pedido & \cod{ORDER\#\allowbreak<order\_id>} & \cod{ITEM\#<nnn>} & \cod{product\_id}, \cod{seller\_id}, preço, frete, categoria do produto \\ \hline
Pagamento & \cod{ORDER\#\allowbreak<order\_id>} & \cod{PAYMENT\#\allowbreak<nnn>} & tipo, parcelas, valor, \cod{GSI2PK}, \cod{GSI2SK} \\ \hline
Avaliação & \cod{ORDER\#\allowbreak<order\_id>} & \cod{REVIEW\#\allowbreak<review\_id>} & nota, comentário, datas \\ \hline
Cliente & \cod{CUSTOMER\#\allowbreak<customer\_id>} & \cod{PROFILE} & \cod{customer\_unique\_id}, CEP, cidade, UF \\ \hline
Produto & \cod{PRODUCT\#\allowbreak<product\_id>} & \cod{PROFILE} & categoria, dimensões, peso \\ \hline
Vendedor & \cod{SELLER\#\allowbreak<seller\_id>} & \cod{PROFILE} & CEP, cidade, UF \\ \hline
\hspace{0pt}Geolocalização & \cod{GEO\#\allowbreak<zip\_prefix>} & \cod{<sequencial>} & latitude, longitude, cidade, UF \\ \hline
\end{tabularx}
\end{table}
 
O pedido funciona como agregado central do modelo: o item de metadados, os
itens comprados, os pagamentos e a avaliação compartilham a mesma chave de
partição (\texttt{ORDER\#<order\_id>}) e ficam fisicamente armazenados
juntos. Assim, uma única operação \texttt{Query} recupera todos os dados de
um pedido, o que no modelo relacional exige consultar quatro tabelas. Os
números sequenciais nas chaves de ordenação foram preenchidos com zeros à
esquerda (por exemplo, \texttt{ITEM\#001}) para que a ordenação
lexicográfica coincida com a numérica.
 
Como o DynamoDB não realiza junções, alguns atributos foram
desnormalizados: a cidade e a UF do cliente foram copiadas para o item de
metadados do pedido, e a categoria do produto foi copiada para cada item do
pedido. Isso evita leituras adicionais no acesso por chave, ao custo de
redundância e de possível inconsistência caso esses atributos sejam
alterados, cenário considerado improvável para dados históricos de vendas.
Clientes, produtos, vendedores e geolocalização foram mantidos como itens
próprios para que a carga do dataset completo seja equivalente nas duas
tecnologias, embora nenhum padrão de acesso avaliado dependa deles
diretamente.
 
Quanto aos tipos, datas foram armazenadas como \textit{strings} no formato
ISO~8601 (\texttt{AAAA-MM-DD hh:mm:ss}), que preserva a ordem cronológica na
comparação lexicográfica e permite consultas por intervalo na chave de
ordenação, e valores financeiros foram armazenados com o tipo
\textit{Number}. Todos os itens ficam muito abaixo do limite de 400~KB. 

Para atender às consultas que não partem do identificador do pedido, foram
criados dois índices secundários globais:
 
\begin{itemize}
\item \textbf{GSI1 (entregas por mês):} chave de partição
\texttt{GSI1PK = DELIVERED\#<AAAA-MM>} e chave de ordenação
\texttt{GSI1SK = order\_delivered\_customer\_date}. Os atributos só são
preenchidos em pedidos com data de entrega, o que torna o índice esparso: pedidos não entregues não ocupam espaço nem consomem
capacidade nele. A consulta de pedidos entregues em janeiro de 2018 se
reduz a uma \texttt{Query} com \texttt{GSI1PK = DELIVERED\#2018-01}.
A granularidade mensal foi escolhida para distribuir os pedidos entre
várias partições, evitando concentrar todo o índice em uma única chave.
Assim como o GSI2, o índice utiliza a projeção \texttt{KEYS\_ONLY}.
 
\item \textbf{GSI2 (pagamentos por tipo):} chave de partição
\texttt{GSI2PK = PAYTYPE\#<payment\_type>} e chave de ordenação
\texttt{GSI2SK = payment\_value}, com projeção \texttt{KEYS\_ONLY}. Como o
valor pago é a própria chave de ordenação, o índice contém tudo o que a
agregação precisa, e cada item projetado é pequeno, reduzindo o consumo de
capacidade de leitura. Os valores mínimo e máximo de cada tipo podem ser
obtidos com uma \texttt{Query} limitada a um item em ordem crescente ou
decrescente; a média, por outro lado, exige ler todas as entradas do tipo.
\end{itemize}
 
 
A consulta de agregação evidencia a principal limitação do modelo: o
DynamoDB não oferece operações como \texttt{GROUP BY}, \texttt{AVG},
\texttt{MIN} ou \texttt{MAX}. A agregação foi, portanto, realizada na
aplicação, com uma \texttt{Query} no GSI2 para cada forma de pagamento,
seguida do cálculo das estatísticas e da ordenação pela média no cliente.
 
 
Por fim, a tabela do DynamoDB foi configurada no modo de capacidade
\textit{on-demand}, no qual a AWS ajusta automaticamente a vazão conforme a
carga, dispensando o dimensionamento prévio de unidades de leitura e
escrita e evitando que limites de capacidade provisionada interfiram nas
medições.

\subsection{Operações por Questão}

A Tabela~\ref{tab:padroes} resume como cada questão comparativa é atendida em
cada tecnologia.

\begin{table}[htbp]
\centering
\caption{Operações utilizadas em cada questão comparativa}
\label{tab:padroes}
\footnotesize
\begin{tabularx}{\textwidth}{|>{\raggedright\arraybackslash}p{2.7cm}|L|L|}
\hline
\textbf{Questão} & \textbf{MySQL} & \textbf{DynamoDB} \\ \hline
Acesso por chave & \cod{SELECT} em \cod{orders}, \cod{order\_items}, \cod{order\_payments} e \cod{order\_reviews} filtrando por \cod{order\_id} & Uma \cod{Query} com \cod{PK = ORDER\#<id>} \\ \hline
Inserção de pedido & \cod{INSERT} em \cod{orders}, \cod{order\_items} e \cod{order\_payments} dentro de uma transação & \cod{TransactWriteItems} com os itens de metadados, itens do pedido e pagamento \\ \hline
Atualização de status & \cod{UPDATE orders SET order\_status ... WHERE order\_id} & \cod{UpdateItem} em \cod{PK = ORDER\#<id>}, \cod{SK = METADATA} \\ \hline
Carga do dataset & \cod{LOAD DATA LOCAL INFILE} por tabela & \cod{BatchWriteItem} em lotes de 25 itens \\ \hline
Entregues em jan/2018 & \cod{SELECT} com intervalo em \cod{order\_delivered\_\allowbreak customer\_date} (uso de índice) & \cod{Query} no GSI1 com \cod{GSI1PK = DELIVERED\#\allowbreak 2018-01} \\ \hline
Estatísticas por forma de pagamento & \cod{SELECT} com \cod{AVG}, \cod{MIN}, \cod{MAX}, \cod{GROUP BY payment\_type} e \cod{ORDER BY} média decrescente & Uma \cod{Query} no GSI2 por tipo de pagamento, com agregação e ordenação na aplicação \\ \hline
Busca de texto (1 tabela) & \cod{SELECT} em \cod{order\_reviews} com \cod{LIKE '\%atraso\%'} & \cod{Scan} com filtro \cod{contains} no comentário \\ \hline
Busca de texto (2 tabelas) & \cod{JOIN} de \cod{orders} e \cod{customers} com \cod{LIKE '\%sao\%'} na cidade & \cod{Scan} com filtro \cod{contains} na cidade, desnormalizada no item do pedido \\ \hline
Agregação por UF (3 tabelas) & \cod{JOIN} de \cod{order\_items}, \cod{orders} e \cod{customers}, com \cod{AVG}, \cod{MIN}, \cod{MAX} e \cod{GROUP BY customer\_state} & \cod{Scan} dos itens e dos pedidos, com junção, agregação e ordenação na aplicação \\ \hline
\end{tabularx}
\end{table}

\clearpage
\section{Resultados Detalhados}
\label{anexo:resultados}

Resultados da execução preliminar descrita na Seção~4.3. Em todas as questões,
os dois bancos retornaram o mesmo número de linhas.

\begin{table}[htbp]
\centering
\caption{Tempo de carga por entidade (s)}
\label{tab:carga}
\small
\begin{tabular}{|l|r|r|r|}
\hline
\textbf{Entidade} & \textbf{Registros} & \textbf{MySQL} & \textbf{DynamoDB} \\ \hline
Clientes   & 99.441  & 1,79  & 95,79  \\ \hline
Produtos   & 32.951  & 0,51  & 32,16  \\ \hline
Vendedores & 3.095   & 0,17  & 4,06   \\ \hline
Pedidos    & 99.441  & 2,33  & 102,65 \\ \hline
Itens      & 112.650 & 2,15  & 109,12 \\ \hline
Pagamentos & 103.886 & 1,49  & 98,83  \\ \hline
Avaliações & 99.224  & 1,86  & 97,60  \\ \hline
\textbf{Total} & \textbf{550.688} & \textbf{10,30} & \textbf{540,23} \\ \hline
\end{tabular}
\end{table}

\begin{table}[htbp]
\centering
\caption{Tempo médio das consultas e número de linhas retornadas (ms, média de 3 execuções)}
\label{tab:consultas-completa}
\small
\begin{tabular}{|l|r|r|r|r|}
\hline
\textbf{Questão} & \textbf{Linhas} & \textbf{MySQL} & \textbf{DynamoDB} & \textbf{Razão} \\ \hline
Acesso por chave                  & 4      & 9,25   & 599,70    & 65$\times$  \\ \hline
Inserção de pedido                & 4      & 17,52  & 174,83    & 10$\times$  \\ \hline
Atualização de status             & 1      & 11,26  & 161,45    & 14$\times$  \\ \hline
Entregues em jan/2018             & 6.597  & 43,71  & 1.541,95  & 35$\times$  \\ \hline
Estatísticas por pagamento        & 5      & 45,19  & 22.441,57 & 497$\times$ \\ \hline
Busca de texto (1 tabela)         & 305    & 69,08  & 23.044,77 & 334$\times$ \\ \hline
Busca de texto (2 tabelas)        & 21.092 & 295,94 & 23.674,01 & 80$\times$  \\ \hline
Agregação por UF (3 tabelas)      & 27     & 775,44 & 40.423,14 & 52$\times$  \\ \hline
\end{tabular}
\end{table}

\end{document}
